\chapter{Redes Neuronales}
\label{cap:Redes Neuronales}

Este capítulo tiene como objetivo establecer los cimientos 
técnicos y la notación que guiarán el resto del trabajo. 
Comenzaremos formalizando la notación básica de las 
\textbf{redes neuronales}, desde sus estructuras más simples 
hasta las arquitecturas profundas, para comprender cómo los 
datos se transforman. Posteriormente, definiremos varias funciones
que nos permitirán entender la arquitectura del transformador en
los siguientes capítulos.
\section{Redes Neuronales}

Consideremos una red neuronal que mapea un vector de entrada
$\mathbf{x} \in \mathbb{R}^{D_i}$  %%vector de entrada
 a un vector de salida 
$\mathbf{y} \in \mathbb{R}^{D_o}$, 
expresada como 
$ \y = f[\x,\phi]$, 
donde $\phi$ representa el conjunto de parámetros del modelo\footnote{ Para evitar confusiones, cabe señalar que los 
subíndices $i$ y $o$ en las variables $D_i$ y $D_o$ provienen de 
los términos en inglés 
\textit{input} y \textit{ouput}, haciendo referencia a la 
dimensión de entrada y dimensión de salida, respectivamente.}. 
La arquitectura se define por la presencia de $K$ capas ocultas, 
cada una de ellas con una dimensionalidad $D_k$ para 
$k = 1,\dots, K$ 
que representa el número de unidades ocultas o neuronas.

Los elementos que componen la red son:
\begin{itemize}
    \item \textbf{Vectores de activación}: $\mathbf{h}_k \in D_k$ representa el estado de las unidades ocultas en la capa $k$.

    \item \textbf{Matrices de peso}: $\boldsymbol{\Omega}_k$ describe la transformación lineal entre la capa $k$ y la $k+1$.

    \item \textbf{Vectores de sesgo (\textit{bias})}: $\boldsymbol{\beta}_k$ representa el término de compensación de cada capa.

    \item  \textbf{Función de activación}: $a[\cdot]$ es una funcion no lineal aplicada elemento a elemento.
\end{itemize}


El flujo de propagación de la información a través de la red se 
define mediante la siguiente secuencia de transformaciones:
$$\begin{aligned}
    \mathbf{h}_1 &= a[\boldsymbol{\beta}_0 + \boldsymbol{\Omega}_0 \mathbf{x}] \\
    \mathbf{h}_2 &= a[\boldsymbol{\beta}_1 + \boldsymbol{\Omega}_1 \mathbf{h}_1] \\
    &\vdots \\
    \mathbf{h}_K &= a[\boldsymbol{\beta}_{K-1} + \boldsymbol{\Omega}_{K-1} \mathbf{h}_{K-1}] \\
    \mathbf{y} &= \boldsymbol{\beta}_K + \boldsymbol{\Omega}_K \mathbf{h}_K
\end{aligned}$$
En la capa $k$-ésima el vector de sesgo $\boldsymbol{\beta}_{k}$ 
tiene dimensión $D_k$ para   \mbox{$k = 1 ,\dots, K-1$,} 
mientras que $\boldsymbol{\beta}_{K}$ tiene dimensión $D_o$. 
La primera matriz de pesos $\boldsymbol{\Omega}_0$ tiene dimensión 
$D_1 \times D_i$ y la última matriz de pesos tiene dimensión $D_o \times D_K$, 
mientras que en el resto de capas intermedias la matriz de pesos 
$\boldsymbol{\Omega}_k$
tiene dimensión $D_{k+1} \times D_k$ (ver figura \ref{fig:red_neuronal}).

%FIGURA RED NEURONAL
\begin{figure}[t]
    \vspace{0.8cm}
    \centering
    \resizebox{\textwidth}{!}{%
        \shorthandoff{>} % desactiva el carácter activo >

\newcommand{\inputnum}{3} 
 
% Hidden layer neurons'number
\newcommand{\hiddennumA}{4}
\newcommand{\hiddennumB}{2}
\newcommand{\hiddennumC}{3}  
 
% Output layer neurons'number
\newcommand{\outputnum}{2} 

 
\begin{tikzpicture}
 
% Input Layer
\foreach \i in {1,...,\inputnum}
{
    \node[circle, 
        minimum size = 8mm,
        fill=orange!30] (Input-\i) at (0,-\i) {$x_\i$};
}

\node at (0,-5) {Entrada, $\x$}; 
\node at (0,-5.5) {$D_i = 3$}; 
\node at (1.8,-4) {$\boldsymbol{\Omega}_0 \in \mathbb{R}^{4 \times 3}$}; 






 
% Hidden Layer 1
\foreach \i in {1,...,\hiddennumA}
{
    \node[circle, 
        minimum size = 8mm,
        fill=teal!50,
        yshift=(\hiddennumA-\inputnum)*5 mm
    ] (HiddenA-\i) at (3.5,-\i) {};
}

\node at (3.5,-5) {Capa Oculta, $\mathbf{h}_1$}; 
\node at (3.5,-5.5) {$D_1 = 3$}; 
\node at (5.3,-4) {$\boldsymbol{\Omega}_1 \in \mathbb{R}^{2 \times 4}$}; 





% Hidden Layer 2
\foreach \i in {1,...,\hiddennumB}
{
    \node[circle, 
        minimum size = 8mm,
        fill=teal!50,
        yshift=(\hiddennumB-\inputnum)*5 mm
    ] (HiddenB-\i) at (7,-\i) {};
}

\node at (7,-5) {Capa Oculta, $\mathbf{h}_2$}; 
\node at (7,-5.5) {$D_2 = 3$};
\node at (8.8,-4) {$\boldsymbol{\Omega}_2 \in \mathbb{R}^{3 \times 2}$}; 


% Hidden Layer 3
\foreach \i in {1,...,\hiddennumC}
{
    \node[circle, 
        minimum size = 8mm,
        fill=teal!50,
        yshift=(\hiddennumC-\inputnum)*5 mm
    ] (HiddenC-\i) at (10.5,-\i) {};
}

\node at (10.5,-5) {Capa Oculta, $\mathbf{h}_3$}; 
\node at (10.5,-5.5) {$D_2 = 3$}; 
\node at (12.3,-4) {$\boldsymbol{\Omega}_3 \in \mathbb{R}^{2 \times 3}$}; 

% Output Layer
\foreach \i in {1,...,\outputnum}
{
    \node[circle, 
        minimum size = 8mm,
        fill=purple!50,
        yshift=(\outputnum-\inputnum)*5 mm
    ] (Output-\i) at (14,-\i) {$y_\i$};
}

\node at (14,-5) {Salida, $\y$}; 
\node at (14,-5.5) {$D_o = 2$};

%COONEXION ENTRADA -> CAPA1
\foreach \i in {1,...,\inputnum}
{
    \foreach \j in {1,...,\hiddennumA}
    {
        \draw[->, shorten >=1pt] (Input-\i) -- (HiddenA-\j);   
    }
}


%COONEXION CAPA1 -> CAPA2
\foreach \i in {1,...,\hiddennumA}
{
    \foreach \j in {1,...,\hiddennumB}
    {
        \draw[->, shorten >=1pt] (HiddenA-\i) -- (HiddenB-\j);   
    }
}


%COONEXION CAPA2 -> CAPA3
\foreach \i in {1,...,\hiddennumB}
{
    \foreach \j in {1,...,\hiddennumC}
    {
        \draw[->, shorten >=1pt] (HiddenB-\i) -- (HiddenC-\j);   
    }
}

%COONEXION CAPA3 -> CAPASALIDA
%COONEXION CAPA2 -> CAPA3
\foreach \i in {1,...,\hiddennumC}
{
    \foreach \j in {1,...,\outputnum}
    {
        \draw[->, shorten >=1pt] (HiddenC-\i) -- (Output-\j);   
    }
}

\end{tikzpicture} 


 % inserta el dibujo
    }
    \caption{Notación para una red con una entrada $\mathbf{x}$ de dimensión $D_i = 3$, 
    una salida $\mathbf{y}$ de dimensión $D_o = 2$ y $K =3$ capas 
    ocultas $ \mathbf{h}_1$, $ \mathbf{h}_2$ y $ \mathbf{h}_3$ de 
    dimensiones $D_1 = 4$, $D_2 = 2$ y $D_3 = 3$ respectivamente.
    Los pesos se almacenan en las matrices $\boldsymbol{\Omega}_k$ que multiplican las 
    activaciones de la capa precedente para generar las pre-activaciones de la siguiente 
    capa. Por ejemplo, la matriz de pesos $\boldsymbol{\Omega}_1$ que computa las 
    pre-activaciones en $\mathbf{h}_2$ a partir de las de $\mathbf{h}_1$ tiene 
    dimensión $2 \times 4$. Esta matriz se aplica a las 4 unidades ocultas de la 
    primera capa y genera las entradas para las 2 unidades de la siguiente. 
    Los sesgos se almacenan en los vectores $\boldsymbol{\beta}_{k}$ y tienen la 
    dimensión de la capa que alimentan.}
    \label{fig:red_neuronal}
\end{figure}

Cuando hablamos de aprendizaje o entrenamiento de un modelo, nos referimos al intento 
de encontrar los parámetros $\phi$ óptimos que permitan realizar predicciones de salida 
precisas y coherentes a partir de los datos de entrada. 
Para aprender estos parámetros, se utiliza un 
conjunto de datos de entrenamiento compuesto por pares de 
ejemplos de entrada y salida $ \{(\mathbf{x}_n, \mathbf{y}_n)\}$. 
Nuestro objetivo es seleccionar aquellos parámetros que mapeen cada entrada de 
entrenamiento a su salida asociada de la forma más precisa posible. 
Para cuantificar el error introducimos la función de pérdida 
$L[\phi]$\footnote{No existe una única función de péridida $L[\phi]$ sino que se define 
según el problema a resolver.}. 
En consecuencia, al entrenar el modelo buscamos el conjunto de parámetros 
$\hat{\phi}$ que minimicen dicha función de pérdida, es decir:
\begin{equation}
    \hat{\phi} = \arg\min_{\phi} \Big[ L[\phi] \Big]
\end{equation}
El operador $\arg \min$ nos indica que el resultado del problema de 
optimización no es el valor mínimo de la pérdida, sino el argumento $\phi$ 
con el cual se alcanza dicho valor mínimo. 



\section{Sotfmax}

La función Sotfmax es una generalización de la función logística.
Dado un vector $\x^T = (\x_1, \dots, \x_n)$, la función Sotfmax $\sigma: \mathbb{R}^n \to (0,1)^n$ se define por componentes como:
\begin{equation}
    \sigma[\x]_i = \frac{e^{\x_i}}{\sum_{j=1}^{n} e^{\x_j}} \quad \text{para } i = 1, \dots, n
\end{equation}
y garantiza que  $\sum_{i=1}^n \sigma[\x]_i = 1$.

Por lo general, y en el campo de la informática, la función Sotfmax también se aplica a matrices $\mathbf{X}  \in \mathbb{R}^{N\times M}$. 
En este contexto, la función se aplica por filas, donde para cada elemento de la matriz se define como:
\begin{equation}
    \sigma[\mathbf{X}]_{i,j} = \frac{e^{\mathbf{X}_{i,j}}}{\sum_{k=1}^{M} e^{\mathbf{X}_{i,k}}}
\end{equation}


\section{Normalización de Capa (\textit{LayerNorm})}
Dado un vector $\mathbf{x}^T = (\x_1, \dots, \x_n) \in \mathbb{R}^n$, la función $\text{LayerNorm}: \mathbb{R}^n \to \mathbb{R}^n$ se define como:
\begin{equation}
\text{LayerNorm}[\mathbf{x}] = \boldsymbol{\gamma} \odot \frac{\mathbf{x} - \mu \mathbf{1}}{\sqrt{\sigma^2 + \epsilon}} + \boldsymbol{\beta}
\end{equation}
donde $\mathbf{1} \in \mathbb{R}^n$ es el vector de unos, $\epsilon > 0$, $\odot$ 
denota el producto de Hadamard (multiplicación componente a componente), y 
$\boldsymbol{\gamma}, \boldsymbol{\beta} \in \mathbb{R}^n$ son vectores de parámetros 
aprendibles. Además:
\begin{equation}\mu = \frac{1}{n} \sum_{j=1}^n \x_j \quad \text{y} \quad \sigma^2 = \frac{1}{n} \sum_{j=1}^n (\x_j - \mu)^2
\end{equation}
La normalización de capa es un operador fundamental diseñado para 
garantizar la estabilidad numérica de los datos a medida que atraviesan las múltiples 
capas de una red neuronal pues, en arquitecturas muy profundas, las multiplicaciones matriciales sucesivas pueden provocar que los valores de 
los vectores crezcan descontroladamente o tiendan a cero. La función LayerNorm mitiga este problema forzando que cada vector adquiera una media 
de cero y una varianza unitaria. Además, una de sus principales diferencias respecto a otras técnicas tradicionales basadas en lotes 
(\textit{batches}), como \textit{Batch Normalization}, es que se aplica de manera independiente sobre cada \textit{token}. Sin embargo, estandarizar estrictamente los datos podría limitar la capacidad expresiva de la red neuronal por lo que, para evitar esta pérdida de flexibilidad, se introducen parámetros que permiten reescalar y desplazar los valores si, durante el entrenamiento, el modelo descubre que una distribución diferente es más óptima para representar la información.










 
 
















